% Part: normal-modal-logic
% Chapter: frame-correspondence
% Section: frames

\documentclass[../../../include/open-logic-section]{subfiles}

\begin{document}

\olfileid{nml}{frd}{fra}

\olsection{चौकटी}

\begin{defn}
  \emph{चौकट} म्हणजे $\mModel{F} = \tuple{W,R}$ ही जोडी, जिथे
  $W$ हा जगांचा रिक्त नसलेला संच आणि $R$ हा~$W$ वरील
  द्विस्थानी संबंध आहे. प्रतिमान~$\mModel{M}$ हे
  $\mModel{F} = \tuple{W,R}$ या चौकटीवर \emph{आधारित}
  आहे तेव्हा आणि केवळ तेव्हाच, जेव्हा काही मूल्यांकन
  $V$ साठी $\mModel{M} = \tuple{W, R, V}$.
\end{defn}

\begin{defn}
  $\mModel{F}$ ही चौकट असेल, तर~$\mModel{F}$ वर आधारित
  प्रत्येक प्रतिमान~$\mModel{M}$ साठी $\mSat{M}{!A}$
  असल्यास $!A$ हे \emph{$\mModel{F}$ मध्ये वैध,}
  म्हणजे $\mModel{F} \Entails !A$, असे म्हणतो.

  $\mClass{F}$ हा चौकटींचा वर्ग असेल, तर त्या वर्गातील
  प्रत्येक चौकट~$\mModel{F} \in \mClass{F}$ साठी
  $\mModel{F} \Entails !A$ असल्यास, आणि तेव्हाच,
  $!A$ हे \emph{$\mClass{F}$ मध्ये वैध,}
  म्हणजे $\mClass{F} \Entails !A$, असे म्हणतो.
\end{defn}

चौकटी महत्त्वाच्या आहेत, कारण रूपबंध आणि प्राप्यता
संबंध~$R$ याचे गुणधर्म यांमधील अनुरूपता \emph{प्रतिमानांच्या
नव्हे, तर चौकटींच्या} पातळीवर असते. उदाहरणार्थ,
\Ax{T} सर्व परावर्ती प्रतिमानांत सत्य असला,
तरी \Ax{T} सत्य असलेल्या प्रत्येक प्रतिमानाचा संबंध
परावर्ती असतोच असे नाही. परंतु असे \emph{खरे} आहे की,
\Ax{T} सर्व परावर्ती \emph{चौकटींवर} \emph{वैध} आहे,
आणि \Ax{T} वैध
असलेली प्रत्येक चौकट परावर्ती आहे, ही दोन्ही विधाने
खरी आहेत.

\begin{rem}
चौकटींच्या वर्गातील वैधता ही प्रतिमानांच्या वर्गातील
वैधतेच्या संकल्पनेचे विशेष प्रकरण आहे: $\mClass{F} \Entails !A$
तेव्हा आणि केवळ तेव्हाच, जेव्हा $\mClass{C} \Entails !A$,
जिथे $\mClass{C}$ हा~$\mClass{F}$ मधील एखाद्या
चौकटीवर आधारित सर्व प्रतिमानांचा वर्ग आहे.

स्पष्टच आहे की, !!a{formula} किंवा रूपबंध वैध असेल,
म्हणजे \emph{सर्व} प्रतिमानांच्या वर्गाच्या संदर्भात वैध
असेल, तर तो चौकटींच्या कोणत्याही वर्ग~$\mClass{F}$
च्या संदर्भातही वैध असतो.
\end{rem}

\end{document}
